\section{Conclusion}
\label{sec:conclusion}

The results show that no single database dominates across all dimensions: IoTDB leads in
sequential write throughput and value-filter queries, TimescaleDB offers constant-latency
downsampling and full SQL compatibility, and InfluxDB occupies a practical middle ground
between performance and memory consumption. The architectural choices of each system, not
configuration tuning, determine scaling behavior, and resource isolation during benchmarking
critically affects result validity.

\subsection{Limitations}
\label{sec:limitations}

The experiments were conducted on a single node, which does not capture the real network
overhead of distributed deployments. Read tests share page-cache state across consecutive
workloads, which may favor later tests. IoTDB's write throughput reflects in-memory
acknowledgement before disk persistence; strict durability configurations may reduce that
figure. Finally, JVM warm-up causes short benchmarks to understate IoTDB's steady-state
throughput.

\subsection{Future Work}
\label{sec:future-work}

Relevant directions include evaluating distributed deployments (multi-node IoTDB, InfluxDB
Clustered, TimescaleDB with Citus), testing with real-world industrial datasets, and
reassessing trade-offs after significant architectural updates such as InfluxDB~3.0's new
columnar engine. Measuring replication and failover costs would also be essential for
production deployment decisions.
